% ============================================================================
% PLANTILLA 04 — PRÁCTICA CALIFICADA
% ----------------------------------------------------------------------------
% Para: prácticas calificadas con ejercicios aplicados, cálculos e
% interpretación de datos y tablas estadísticas.
% Compilar: latexmk -pdf plantilla-04-practica-calificada.tex
% Solucionario: \documentclass[soluciones]{evaluacion}
% ============================================================================
\documentclass{evaluacion}

\universidad{Universidad Nacional de San Cristóbal de Huamanga}
\facultad{Facultad de Ciencias Económicas, Administrativas y Contables}
\escuela{Escuela Profesional de Economía}
\curso{Estadística para Economistas I}
\codigocurso{ES-251}
\docente{Mg. Nombre Apellido}
\periodo{Semestre 2026-I}
\tipoevaluacion{Práctica Calificada 02}
\fechaevaluacion{21 de julio de 2026}
\duracion{100 minutos}
\modalidad{Presencial · individual}

\begin{document}
\membrete

\begin{instrucciones}
\begin{itemize}
  \item Presente el procedimiento completo; los resultados sin desarrollo
        no otorgan puntaje.
  \item Redondee a cuatro decimales. Use la tabla normal estándar adjunta
        al cuadernillo.
  \item La práctica suma \totalpuntos{} puntos.
\end{itemize}
\end{instrucciones}


\begin{ejercicio}[6][Estadística descriptiva]
La siguiente tabla resume el gasto mensual en alimentación (en soles) de
una muestra de 80 hogares de Ayacucho:

\begin{center}
\small
\begin{tabular}{@{}c S[table-format=2.0] S[table-format=3.1]@{}}
  \toprule
  \textbf{Gasto (S/)} & {\textbf{Hogares} ($n_i$)} & {\textbf{Marca} ($x_i$)}\\
  \midrule
  {}[200, 400)  & 12 & 300.0\\
  {}[400, 600)  & 26 & 500.0\\
  {}[600, 800)  & 24 & 700.0\\
  {}[800, 1000) & 13 & 900.0\\
  {}[1000, 1200)&  5 & 1100.0\\
  \bottomrule
\end{tabular}
\end{center}

\begin{apartados}
  \item Calcule la media, la mediana y la desviación estándar muestral. \pts{3}
  \item Calcule el coeficiente de variación e interprete. \pts{1{,}5}
  \item ¿La distribución es simétrica? Justifique con el coeficiente de
        asimetría de Pearson. \pts{1{,}5}
\end{apartados}
\espaciorespuesta{7cm}
\begin{solucion}
\textbf{a)} Media: $\bar{x} = \frac{1}{80}\sum n_i x_i =
\frac{51\,600}{80} = 645$. Mediana: la posición 40 cae en el intervalo
$[600,800)$:
\[
  Me = 600 + \frac{40-38}{24}\,(200) = 616{,}67.
\]
Varianza muestral: $s^{2} = \frac{\sum n_i x_i^{2} - n\bar{x}^{2}}{n-1}
= \frac{37\,860\,000 - 33\,282\,000}{79} = 57\,949{,}37$, de donde
$s = 240{,}73$.

\textbf{b)} $CV = s/\bar{x} = 240{,}73/645 = 0{,}3732$ (37,3\,\%):
dispersión relativa moderada-alta; la media sigue siendo representativa
pero con heterogeneidad apreciable entre hogares.

\textbf{c)} $A_p = 3(\bar{x}-Me)/s = 3(645-616{,}67)/240{,}73 = 0{,}353>0$:
asimetría positiva leve (cola derecha, gastos altos poco frecuentes).
\end{solucion}
\end{ejercicio}

\begin{ejercicio}[7][Probabilidad]
En cierta ciudad, el 60\,\% de los hogares tiene acceso a internet fijo.
Se seleccionan 10 hogares al azar.
\begin{apartados}
  \item Calcule la probabilidad de que exactamente 7 tengan acceso. \pts{2}
  \item Calcule la probabilidad de que al menos 8 tengan acceso. \pts{2}
  \item Si se seleccionaran 400 hogares, aproxime mediante la normal la
        probabilidad de que más de 250 tengan acceso (use corrección por
        continuidad). \pts{3}
\end{apartados}
\espaciorespuesta{7cm}
\begin{solucion}
Sea $X\sim B(10;\,0{,}6)$.

\textbf{a)} $P(X=7) = \binom{10}{7}(0{,}6)^{7}(0{,}4)^{3} = 0{,}2150$.

\textbf{b)} $P(X\ge 8) = P(8)+P(9)+P(10) =
0{,}1209 + 0{,}0403 + 0{,}0060 = 0{,}1673$.

\textbf{c)} Con $n=400$: $\mu = np = 240$,
$\sigma = \sqrt{np(1-p)} = \sqrt{96} = 9{,}7980$. Entonces
\[
  P(X > 250) \approx P\!\left(Z > \frac{250{,}5 - 240}{9{,}7980}\right)
  = P(Z > 1{,}0716) = 0{,}1419.
\]
\end{solucion}
\end{ejercicio}

\begin{ejercicio}[7][Inferencia]
Una muestra aleatoria de 36 microempresas arroja una utilidad media
mensual de S/~2\,450 con desviación estándar muestral de S/~540.
\begin{apartados}
  \item Construya un intervalo de confianza del 95\,\% para la utilidad
        media poblacional. \pts{3}
  \item Contraste al 5\,\% la hipótesis $H_0\colon \mu = 2\,600$ contra
        $H_1\colon \mu < 2\,600$. Concluya en términos del problema. \pts{4}
\end{apartados}
\espaciorespuesta{7cm}
\begin{solucion}
\textbf{a)} Con $t_{0{,}975;\,35} = 2{,}0301$:
\[
  \bar{x} \pm t\,\frac{s}{\sqrt{n}}
  = 2\,450 \pm 2{,}0301\cdot\frac{540}{6}
  = 2\,450 \pm 182{,}71
  \;\Rightarrow\; IC_{95\%} = [2\,267{,}29;\; 2\,632{,}71].
\]
\textbf{b)} Estadístico de prueba:
\[
  t = \frac{2\,450 - 2\,600}{540/\sqrt{36}} = \frac{-150}{90} = -1{,}6667,
  \qquad t_{crit} = -t_{0{,}95;\,35} = -1{,}6896.
\]
Como $-1{,}6667 > -1{,}6896$, no se rechaza $H_0$ al 5\,\%: la evidencia
es insuficiente para afirmar que la utilidad media sea menor que
S/~2\,600 (el valor-$p\approx 0{,}052$ queda apenas por encima de
$\alpha$).
\end{solucion}
\end{ejercicio}

\findelexamen
\end{document}
